\scanpage{4}
\mspage{- 176 -}

n° 280

\[
v(x+y)=\inf\,(v(x),v(y)).
\]

La relation (5) se déduit de $(\mathrm{VA}_{\mathrm{II}})$ par récurrence sur $n$. S'il existe un seul indice $k$ tel que $v(x_k)=\inf_{1\leqslant i\leqslant n}(v(x_i))$, on peut écrire $x_i=y_i\,x_k$ avec $v(y_i)>0$ pour $i\ne k$ ; on en déduit $\sum_{i=1}^n x_i={}
=x_k(1+\sum_{i\ne k}y_k)=x_k(1+y)$ avec $v(y)>0$ ; si $A$ désigne l'anneau de $v$, on a $1+y\in A$ et $1+y\notin\mathfrak{m}(A)$, d'où $v(1+y)=0$ et $v(\sum_{i=1}^n x_i)=v(x_k)$.

\uline{Corollaire} - \uline{Si des éléments} $x_i$ \uline{de} $K$ $(1\leqslant i\leqslant n,\ n\geqslant 2)$ \uline{sont tels que} $\sum_{i=1}^n x_i=0$, \uline{il y a au moins deux indices} $j$ \uline{et} $k$ \uline{tels que} $v(x_j)={}
=v(x_k)=\inf_{1\leqslant i\leqslant n}(v(x_i))$.

S'il n'y avait qu'un seul indice $k$ avec $v(x_k)=\inf_{1\leqslant i\leqslant n}(v(x_i))$, la prop.4 montrerait que $v(x_k)=v(0)=+\infty$, d'où $v(x_i)=+\infty$ pour $1\leqslant i\leqslant n$ contrairement à l'hypothèse faite sur $k$.

2. - \uline{Traductions}

Soient $K$ un corps, $\varphi$ une place de $K$, $v$ une valuation de $K$, et $A$ un anneau de valuation pour $K$. Nous dirons que $A$, $\varphi$ et $v$ sont \uline{associés} si $A$ est l'anneau de $\varphi$ et l'anneau de $v$. En vertu du n°1 et du \S2, chacun des trois objets $A$, $\varphi$ et $v$ détermine alors les deux autres (à une équivalence près en ce qui concerne les places et les valuations)/ On a en particulier
\[
\begin{aligned}
x\in A &\Longleftrightarrow \varphi(x)\ne\infty &\Longleftrightarrow v(x)&\geqslant 0\\
x\in\mathfrak{m}(A) &\Longleftrightarrow \varphi(x)=0 &\Longleftrightarrow v(x)&>0\\
x\in A-\mathfrak{m}(A)&\Longleftrightarrow\varphi(x)\ne0\text{ et }\varphi(x)\ne\infty & v(x)&=0
\end{aligned}
\]

Tout résultat portant sur les anneaux de valuation, les places ou les valuations se traduit en un résultat portant sur les deux autres notions. Ainsi, la prop.4 du \S2 donne :

\uline{Proposition} 5 - \uline{Soient} $K$ \uline{un corps}, $F$ \uline{un sousècorps de} $K$, \uline{et} $v$ \uline{une valuation de} $F$. \uline{Il existe une valuation} $w$ \uline{de} $K$ \uline{dont la restriction à} $F$ \uline{est équivalente à} $v$.

% archive sheet number: 3

